% Cross-chapter wording: the standalone build treats this as Chapter 2;
% thesis-draft.tex redefines these macros when it nests it as Section 1.2.
\providecommand{\unitname}{chapter}
\providecommand{\bioref}{Chapter~1}
\providecommand{\biosec}[2]{#1}
\providecommand{\gapref}{Chapter~3}
\chapter{Computational Foundation}
\label{ch:comp}

\section{Scope and Review Axes}
\label{sec:comp-scope}

\bioref{} described the circuit as a whole: the laminar column, its excitatory
and inhibitory cell classes, and the cortico-thalamo-cortical loop. This
\unitname{} turns to the computational side and asks how existing models represent
cortical inhibition, sparse coding and associative memory, and what they have
accomplished. Rather than cataloguing these models, it reads each one for what
it establishes and also for where it stops, so that the boundary of current
work is visible when the gap is stated.

The review is organised along three axes, the first of which is the
\emph{substrate}: the medium in which a model represents and computes.
Candidates include binary sparse distributed representations over minicolumns
with dendritic segments, rate-coded populations, spiking attractor networks,
and multi-compartment neurons. The substrate matters because a mechanism shown
in one medium does not automatically carry over to another. The second axis is
the \emph{inhibition model}. Inhibition is either homogeneous, a single
mechanism such as k-winners-take-all applied uniformly, or differentiated,
with distinct interneuron classes that differ in target compartment, dynamics
and arithmetic effect. The third axis is \emph{encoding versus retrieval}. It
asks how a model relates the storage of new associations to the recall of
stored ones, and what, if anything, decides between them.

Because the thesis turns on these three terms, they are defined here.
\emph{Encoding} is the process by which a new input pattern is turned into a
stable representation and is stored, here through synaptic change.
\emph{Retrieval} is the process by which a partial or noisy cue reactivates a
stored representation. The two are not mutually exclusive: a network can
retrieve while its synapses remain plastic, and encoding a new input usually
involves matching it against what is already stored. \emph{Mode arbitration}
is the mechanism that decides which of the two dominates at a given moment,
for example by gating plasticity or by weighting external input against
recurrent or feedback input. The third axis therefore asks less whether a
model encodes and retrieves than whether it arbitrates between the two, and by
what means.

The third axis has a clear precedent outside the neocortex. In the hippocampal
model of Hasselmo et al.~\cite{hasselmo2002theta}, the theta rhythm
alternates between two phases. In one, entorhinal input is strong and
long-term potentiation is strong, but recurrent input from CA3 is weak, so new
associations are encoded without interference from old ones. In the other, the
balance reverses and stored associations are retrieved. That separation is
implemented by oscillatory and cholinergic modulation in the hippocampus.

Cutsuridis et al.~\cite{cutsuridis2010encoding} showed how interneuron classes
could carry out this separation. Their biophysical model of the CA1 microcircuit
contains pyramidal cells and four interneuron types. Basket and axo-axonic
cells, the perisomatic and axon-targeting classes that in CA1 largely express
PV, fire during the storage half of the theta cycle. Bistratified and oriens
lacunosum-moleculare (OLM) cells, the latter expressing SST, fire during the
recall half. During recall, OLM cells inhibit the distal dendrites on which
entorhinal input arrives, protecting the recalled pattern from interfering
input, while bistratified cells set the threshold for recall. Removing
interneuron classes one at a time tested these roles: recall survived the loss
of basket and axo-axonic inhibition, but degraded under interfering entorhinal
input when OLM inhibition was removed. The interneurons divide the work of the
two modes between them, but the phase that selects a mode is still set from
outside the circuit, by septal input pacing the theta rhythm. Whether a
neocortical column could divide encoding and retrieval between its interneuron
classes in a comparable way is the question this review keeps in view.

In the remainder of the \unitname{}, Section~\ref{sec:comp-htm} reviews the HTM
tradition and its single model of inhibition. Section~\ref{sec:comp-assoc}
covers associative extensions of that tradition, Sparse Distributed Memory and
the Neural Association Algorithm. Section~\ref{sec:comp-interneuron} reviews
circuit models that assign distinct roles to PV, SST and VIP interneurons.
Section~\ref{sec:comp-attractor} reviews associative-memory columns in the
attractor and BCPNN tradition. Section~\ref{sec:comp-operators} examines the
two operator mechanisms at issue, divisive normalisation and dendritic
arbitration. Section~\ref{sec:comp-synthesis} places all of these models on
the three axes.

\section{The HTM, CLA and SDR Tradition}
\label{sec:comp-htm}

Hierarchical Temporal Memory (HTM) and its Cortical Learning Algorithms (CLA)
form a family of models built on a small set of neocortical principles: sparse
distributed representations, minicolumns, pyramidal neurons with active
dendrites, and continuous local learning~\cite{hawkins2011cla}. The tradition
matters to this thesis because it provides the substrate that the thesis builds
on. It is equally significant for what it leaves out.

\subsection{Sparse Distributed Representations}
\label{subsec:comp-sdr}

A sparse distributed representation (SDR) is a large binary vector in which
only a small fraction of bits are active at any time, typically a few percent.
Ahmad and Hawkins~\cite{ahmad2015sdr} derived properties that make SDRs
suitable for cortical computation, focusing on scaling, robustness, and
generalisation. The overlap between two SDRs measures their similarity. The
probability that two unrelated SDRs share many active bits falls rapidly as
dimensionality grows, so a pattern can be recognised reliably by sampling only
a small subset of its active bits. The same property lets a set of SDRs be
stored as their union and still be matched with a low false-positive rate.
Ahmad and Hawkins~\cite{ahmad2016sdr} later linked these properties to
dendritic segments, showing that a segment which detects a coincidence among a
small sample of active inputs behaves as a robust pattern detector over SDRs.

\subsection{The Spatial Pooler}
\label{subsec:comp-sp}

The Spatial Pooler (SP) converts arbitrary binary input into an SDR over a set
of minicolumns~\cite{cui2017sp}. Each minicolumn has a pool of potential
feedforward synapses, each with a scalar permanence. A synapse counts as
connected when its permanence exceeds a threshold, and the minicolumn's overlap
is the number of connected synapses whose inputs are active. A minicolumn
becomes active if its overlap exceeds a stimulus threshold and ranks within the
top $s$ percent of its neighbourhood, with $s$ typically 2\%. The authors
describe this selection as a k-winners-take-all computation. When the
neighbourhood covers the whole layer, it becomes global inhibition. Learning is
Hebbian: winning minicolumns strengthen synapses from active inputs and weaken
synapses from inactive ones. A slower homeostatic mechanism, boosting, raises
the excitability of minicolumns that are active too rarely. Cui et al.\ showed
that the Spatial Pooler adapts quickly to changing input statistics, becomes
more robust to noise with learning, uses its minicolumns efficiently, and
degrades gracefully under cell death.

Dobric et al.~\cite{dobric2021homeostatic,dobric2022newborn} found that the
same homeostatic mechanism can destabilise learned representations, pushing the
Spatial Pooler into unstable, ``epileptic'' activity. Their remedy confines
homeostatic plasticity to an initial ``newborn stage'' and then disables it,
after which learned SDRs remain stable.

\subsection{Temporal Memory and Active Dendrites}
\label{subsec:comp-tm}

Hawkins and Ahmad~\cite{hawkins2016neurons} extended the HTM neuron with three
integration zones. The proximal zone receives feedforward input and drives
action potentials. The basal zone receives contextual input, mostly from
nearby cells in the same region, and the apical zone receives feedback. Each
dendritic segment acts as an independent pattern detector, so a single neuron
can recognise hundreds of distinct patterns. A match on a basal or apical
segment does not fire the cell. It depolarises the cell slightly, putting it
in a predictive state.

The Temporal Memory (TM) network uses this predictive state for sequence
learning. All cells in a minicolumn share the feedforward receptive field
selected by the Spatial Pooler. When a minicolumn becomes active, its
depolarised cells fire first and inhibit their neighbours, so the input is
represented by a sparse, context-specific subset of cells. If no cell in the
minicolumn was predicted, all of its cells become active, which signals an
unanticipated input. Because predictions are sparse, the network can hold many
simultaneous predictions without confusion. Cui et al.~\cite{cui2016sequence}
showed that Temporal Memory learns variable-order sequences, keeps multiple
predictions open until disambiguating input arrives, and matches the prediction
accuracy of ARIMA (autoregressive integrated moving average) models, extreme
learning machines, LSTM and echo-state networks on artificial and real
streaming data.

\subsection{Continual Unsupervised Learning}
\label{subsec:comp-continual}

Learning in this lineage is local, Hebbian and continuous. There is no global
error signal and no separation between training and testing. Cui et
al.~\cite{cui2016sequence} reported continuous online learning, robustness to
sensor noise, fault tolerance, and good performance without task-specific
hyperparameter tuning. Iyer et al.~\cite{iyer2022dendrites} carried the
dendritic principle into deep networks for continual learning. They gated units
through context-modulated dendritic segments and enforced k-winners-take-all
sparsity, and the combination reduced catastrophic forgetting across
sequentially learned tasks.

The same property bears on the third review axis, though less directly than it
might seem. Continuous learning does not make encoding and retrieval
inseparable. Temporal Memory retrieves, in the sense of predicting the next
element from stored sequences, while its synapses remain plastic, and each new
input is matched against stored segments as it is learned. Plasticity is
applied on every cycle, and nothing in the network decides that a given input
should be stored rather than only recalled, or the reverse. In the terms of
Section~\ref{sec:comp-scope}, HTM encodes and retrieves, but does not
arbitrate between the two.

\subsection{Inhibition in the HTM Lineage}
\label{subsec:comp-htm-inhibition}

Across the Spatial Pooler, Temporal Memory and their continual-learning
extensions, inhibition is a single operation that ranks the candidates and
keeps the top $k$, applied either globally or within a neighbourhood. It
has no cell identity, no target compartment and no dynamics. Its effect is set
by rank rather than by strength: strict kWTA imposes an upper bound of
$k$ on the active population and selects among candidates by their
ordering, not by how strongly they are driven. Fewer than $k$ units can
be active when eligibility rules intervene, such as the stimulus threshold or
a local neighbourhood with few candidates, but these limits come from the
eligibility criteria, not from the inhibition itself. Inhibition also does not
act separately on the apical zone, even though the HTM neuron has one, so it
offers no means of weighing feedback against feedforward drive.

This is a modelling choice, not a claim about biology, and the authors' own
discussions point to more structure than the algorithms use. Hawkins and
Ahmad~\cite{hawkins2016neurons} propose fast-spiking basket cells as the most
likely source of the rapid inhibition that TM depends on. These are the
parvalbumin-expressing perisomatic cells described in Section~\biosec{1.3}{sec:bio-inhibitory}. Cui et
al.~\cite{cui2017sp} note that the cortex contains many distinct classes of
inhibitory neurons and suggest that classes innervating cells within a
minicolumn, such as bipolar and double bouquet cells, may help define it. The
lineage therefore attributes its inhibition to specific and differing cell
types in biology, but implements a single homogeneous rule in computation.

This contrast with the biological circuit raises two questions for the present
work. The first is whether replacing k-winners-take-all with a PV-like
divisive operator changes how assembly size is set, from an upper bound
imposed by rank to a recruitment that depends on input magnitude, and under
what conditions it can suppress the population entirely. The second, taken up
together with the first in Section~\ref{sec:comp-operators}, is whether adding
an SST-like dendritic operator gives the network a handle on feedback that it
currently lacks.

\section{Associative Extensions}
\label{sec:comp-assoc}

The HTM algorithms reviewed above encode inputs and learn sequences within a
single region. Two lines of work take sparse binary representations further,
from representation to association. Sparse Distributed Memory is an earlier and
independent formalism for storing and recalling patterns in a high-dimensional
binary space. The Neural Association Algorithm extends the HTM tradition itself
to associations across cortical areas.

\subsection{Sparse Distributed Memory}
\label{subsec:comp-sdm}

Kanerva's Sparse Distributed Memory (SDM) stores binary words of very high
dimension, typically around a thousand bits~\cite{kanerva1988sdm}. The address
space is far too large for every address to have physical storage, so the
memory is built from a fixed random sample of addresses, called hard locations,
each holding one counter per bit. A word is written by activating every hard
location within a given Hamming radius of its address and adding the word to
their counters, incrementing for ones and decrementing for zeros. A word is read
by activating the hard locations around the cue address, summing their counters
bit by bit, and thresholding the sums. Because each word is spread across many
locations and each location holds traces of many words, reading is tolerant to
noisy cues. Reading repeatedly, with each output used as the next address,
converges to the stored word when the cue lies close enough to it. Storing each
item at the address of its predecessor turns the same mechanism into a sequence
memory. Kanerva related the architecture to the circuitry of the cerebellum.

SDM matters to this review in three ways. First, its hard locations play a
role similar to the minicolumns of the Spatial Pooler, as both map an input
onto a sparse set of units chosen by similarity. They differ in how that set
is chosen. The SDM addresses are fixed and random rather than learned, and the
active set is every location inside a radius, so its size varies with the cue.
The Spatial Pooler learns its receptive fields but bounds the active set at
$k$ by rank. SDM therefore needs no inhibition to stay sparse, and its
active population grows or shrinks with the cue, a property that rank-based
selection does not provide (Section~\ref{subsec:comp-htm-inhibition}). Second,
SDM separates writing from reading cleanly, as two distinct operations. That
arbitration is imposed from outside, however: the caller decides whether a
cycle writes or reads, and nothing in the memory makes that decision. Third,
SDM is an abstract memory and not a model of cortical tissue. It has no
minicolumns, dendrites or laminae, so it says nothing about how a cortical
circuit might implement either its selection rule or its write and read modes.

A note on terminology is needed here. The title of this thesis describes its
model as a laminar sparse-distributed memory column. The term is used in its
descriptive sense: a cortical column that stores and recalls associations over
sparse distributed representations. The column is built on the HTM substrate
of Section~\ref{sec:comp-htm}, not on Kanerva's architecture of random hard
locations and counters. The two share the principle that makes such memories
robust, namely storage distributed across many sparsely active units.
Kanerva's model is reviewed here because it is the original statement of that
principle, and because it makes explicit two properties that the HTM lineage
lacks: an active set whose size varies with the input, and an explicit choice
between writing and reading.

\subsection{The Neural Association Algorithm}
\label{subsec:comp-naa}

The Neural Association Algorithm (NAA) extends HTM from a single region to a
cortical space of many areas, each running the same
algorithm~\cite{dobric2024naa}. In each area, an encoder, usually the Spatial
Pooler, produces an SDR. An association function then maps the active cells of
any area onto a sparse set of active cells in a target area. NAA uses the
three segment types of the HTM neuron for this. Proximal segments connect to
sensory input, distal (basal) segments connect to cells of other minicolumns
in the same area, and apical segments connect to cells in other areas. Each
segment is active or matching, depending on how many connected synapses it has
relative to two thresholds. NAA distinguishes two kinds of context: a spatial
context is an undirected association between two SDRs, such as a song and the
place where it was heard, whereas a temporal context is a directed
association, created when one state follows another. Learning happens only
inside active minicolumns, where synapses from active cells onto active
segments are strengthened and new synapses are grown, while synapses onto
matching segments in inactive minicolumns are weakened. Alongside the
association model, the work introduces the newborn stage of the Spatial Pooler
(Section~\ref{subsec:comp-sp}) and a parallel implementation of minicolumn
computation based on the actor programming model.

NAA is the closest predecessor of this thesis. It shares the substrate and aims
at association, which an SDR column must support before encoding and retrieval
can be told apart. Its scope on the other two axes is stated openly by its
author. On inhibition, the implementation ``omits using explicit inhibitory
cells for performance reasons'' and replaces them with an inhibition
algorithm~\cite[Sec.~5.7.3]{dobric2024naa}. That algorithm comes in a global and
a local version, both selecting the minicolumns with the highest overlap at a
fixed density, which is the k-winners-take-all rule of
Section~\ref{subsec:comp-htm-inhibition}. No significant difference in learning
performance was found between the versions, and the thesis names the
investigation of inhibition as future work, specifically how the inhibition
process influences sparsity and the overall performance of
learning~\cite[Sec.~9.5.1.7]{dobric2024naa}. On encoding versus retrieval, NAA
learns in every cycle, as HTM does. Recalling an input from its SDR, which the
thesis calls reverse encoding, is currently handled by a classifier component
that is not aligned with biology, and replacing it with learned associations is
again listed as future work~\cite[Sec.~9.5.1.2]{dobric2024naa}. Finally, only
the single-area core of NAA is implemented. The apical segments that would link
areas, and with them any channel for feedback, are part of the model but not
yet of the implementation~\cite[Sec.~9.5.1.1]{dobric2024naa}.

Its homogeneous inhibition, which the author explicitly leaves open, is a
deliberate simplification of a framework built for scale, not an oversight.

\section{Interneuron-Resolved Circuit Models}
\label{sec:comp-interneuron}

Alongside the HTM tradition, a separate body of work models cortical inhibition
at the level of interneuron classes. These models take the connectivity
described in Section~\biosec{1.3}{sec:bio-inhibitory} as their starting point. PV cells inhibit the
perisomatic region of pyramidal cells, SST cells inhibit their dendrites, and
VIP cells inhibit SST cells and so release the dendrites from inhibition. The
models then ask what this division of labour computes. They are reviewed here in
four groups: disinhibitory gating, state-dependent responses to modulation,
learning and optimisation within the interneuron circuit, and large-scale
laminar models.

\paragraph{Disinhibitory gating.}
Yang et al.~\cite{yang2016gating} proposed that the SST and VIP motif gates
input pathways on the dendrites of pyramidal cells. Their pyramidal neuron has a
spiking soma and several independent dendritic compartments, reduced for
network simulations to a multi-compartment rate model. SST cells keep every
dendritic branch inhibited by default, and top-down control signals onto VIP
cells release selected branches. If the inputs of one pathway cluster on the
same branches, disinhibiting those branches gates that pathway on while the
others stay closed. This selectivity survives dense and even random interneuron
connectivity. A calcium-based plasticity rule regulated by dendritic inhibition
produces the required clustering, and the gated circuit performs a
context-dependent decision task in which a rule signal selects whether motion or
colour determines the response.

Hertäg and Sprekeler~\cite{hertag2019amplifying} studied how the same motif
redistributes inhibition between the soma and the dendrite. Their rate model
uses a two-compartment pyramidal cell, with PV cells inhibiting the soma, SST
cells inhibiting the apical dendrite, and VIP cells inhibiting only SST cells.
Mutual inhibition between VIP and SST cells turns weak inputs to VIP cells into
large shifts between somatic and dendritic inhibition. When that mutual
inhibition is strong, the circuit becomes bistable: a transient input switches
it persistently between two processing modes, one in which top-down input to
the apical dendrite is integrated and one in which it is cancelled. Of the
models in this section, this is the closest analogue to the arbitration that
this thesis is concerned with, a circuit-level switch that decides how feedback
is treated. It is framed as a mode of processing top-down input, however, not as
a choice between storing and recalling a memory, and the model learns no
associations.

\paragraph{State-dependent responses.}
Garcia del Molino et al.~\cite{garciadelmolino2017paradoxical} built a rate
model of pyramidal, PV, SST and VIP populations to explain a puzzling
observation from mouse visual cortex. The response of SST cells to top-down
behavioural modulation, which targets VIP cells, can change sign when the
visual input changes. Their analysis traces this response reversal to specific
connection strengths and shows that the effect of activating one interneuron
class depends on the state of the whole network. For this thesis the point is a
cautionary one. A single interneuron class cannot be assigned one fixed effect
from its connectivity alone, so any abstraction of PV or SST as one operator is
a simplification whose limits must be stated.

\paragraph{Learning and optimisation.}
Two models show that the differentiation itself can be learned or derived.
Hertäg and Sprekeler~\cite{hertag2020prediction} added inhibitory plasticity to
a circuit of pyramidal, PV, SST and VIP cells and showed that it learns to
produce prediction-error neurons, which signal mismatches between sensory input
and a prediction. Keijser and Sprekeler~\cite{keijser2022optimizing} took a
normative route. They optimised the connectivity and short-term plasticity of a
spiking interneuron circuit for compartment-specific feedback inhibition, and
the interneurons diversified during optimisation into two classes resembling PV
and SST cells. The resulting circuit acts as a decoder that inverts the
nonlinear computation within the pyramidal cell. This gives computational
support for treating the PV and SST split as a functional one: two classes
emerge when inhibition must be specific to compartments.

\paragraph{Laminar models.}
At the largest scale, Billeh et al.~\cite{billeh2020systematic} integrated
anatomical and physiological data into models of mouse primary visual cortex
spanning all layers, with about 230,000 neurons in 17 cell classes. The model
exists in a biophysically detailed version and a point-neuron version with
identical connectivity, and both were compared against recordings from awake
mice. Moreni et al.~\cite{moreni2025celltype} built a smaller model at the
scale of a single column, drawing its connectivity from the same data. Their
network of 5,000 leaky integrate-and-fire neurons has pyramidal, PV, SST and
VIP populations in layers 2/3, 4, 5 and 6, and VIP cells alone in layer 1.
Comparing two input regimes, they found that feedforward input to layer 4
produced net excitation of the column, whereas feedback-like input to layer 5
produced net inhibition, which the authors attribute to translaminar gain
control through column-wide inhibition from layer 6 (see Section~\biosec{1.4}{sec:bio-canonical}). The
response to perturbing any single cell group also depended on the state of the
column: strong feedforward input made the column less sensitive to every
perturbation, while feedback input modulated interactions within it. How the
column treats its inputs therefore depends on whether it is driven from below
or from above, which bears directly on the arbitration question of this
thesis. Like Billeh et al., however, the model has no synaptic plasticity and
performs no memory task. Laminar models of this kind show that interneuron
detail across layers can be assembled and constrained by data. Their question
is how the circuit responds in different states, not what it stores or
retrieves.

These models provide what the HTM tradition lacks on the inhibition axis. They
assign distinct targets and dynamics to interneuron classes and analyse what
each class contributes. On the other two axes, however, they offer much less.
Their substrate is rate-coded or spiking populations, with no sparse
distributed representations, no minicolumns and no learned assemblies, and
their questions concern gating, modulation, error signalling and physiological
fidelity. None of the models reviewed here stores a set of patterns and then
recalls one from a partial cue.

\section{Associative-Memory Columns}
\label{sec:comp-attractor}

A third line of work asks directly whether networks of realistic cortical
neurons can act as associative memories. It descends from Hebb's cell-assembly
theory and its formal counterpart, the attractor network, in which a stored
pattern is a stable state that the network falls into from a partial or noisy
cue. Since the 1990s, Lansner and colleagues have tested this idea with
increasingly detailed models of cortical layers 2/3, many of them built around a
Bayesian-Hebbian learning rule known as BCPNN (Bayesian Confidence Propagation
Neural Network)~\cite{lansner2009associative}. Because these models store and
recall patterns, they are the only family reviewed so far that addresses the
memory task directly. Their treatment of inhibition has also grown more
differentiated over time.

\paragraph{Columns as memory units.}
Fransén and Lansner~\cite{fransen1998horizontal} modelled the horizontal network
of layers 2/3 as a set of connected cortical columns, with the column rather than
the single neuron as the functional unit. Each column contained biophysically
detailed pyramidal cells and fast-spiking inhibitory interneurons representing
basket cells. A memory was an assembly of columns linked by sparse, long-range
excitatory connections. Connection strengths were not learned during the
simulation. They were computed in advance with the Bayesian rule from the full
set of patterns to be stored. Positive weights became excitatory connections
between the pyramidal cells of different columns. Negative weights were routed
onto the basket cells of the receiving column, so that inhibition between
competing assemblies was disynaptic. With this wiring the network completed an
assembly from a partial cue, resolved rivalry between overlapping assemblies,
and sustained activity after the cue only when the excitability of its pyramidal
cells was raised by simulated neuromodulation. The model thus showed that
biophysically realistic neurons can support attractor memory. It used a single
class of interneuron, and the switch between transient and sustained activity
came from a neuromodulator, not from the inhibitory circuit.

\paragraph{Modular attractor networks.}
Later versions organised the same layer 2/3 network into hypercolumns, each
containing a set of minicolumns~\cite{lundqvist2006attractor}. Basket cells within
a hypercolumn provide lateral inhibition that makes each module behave as a soft
winner-take-all, typically allowing one active minicolumn at a time, and a
stored pattern is a set of minicolumns distributed across hypercolumns.
Regular-spiking non-pyramidal cells, possibly double bouquet cells, provided
inhibition between minicolumns~\cite{krishnamurthy2012cortical}. This
organisation is strikingly close to the minicolumn and k-winners-take-all
structure of the Spatial Pooler (Section~\ref{subsec:comp-sp}). The difference
is that the competition is implemented by a population of basket cells rather
than by a sorting rule.

\paragraph{Differentiated interneurons in memory columns.}
Three studies added further interneuron classes to this architecture.
Krishnamurthy, Silberberg and Lansner~\cite{krishnamurthy2012cortical} introduced
Martinotti cells, which receive strongly facilitating synapses from pyramidal
cells, alongside basket cells, whose input synapses depress. When the network
entered a stored attractor, basket cells fired first and then tapered off, while
Martinotti cells engaged with a delay and shifted the balance of excitation and
inhibition until the attractor terminated. A sensitivity analysis showed
that Martinotti cells, rather than cellular adaptation or synaptic depression,
dominated the attractor dwell time, and replacing them with additional basket
cells did not reproduce the effect. The temporal profile of each interneuron
class, not only the amount of inhibition, therefore shaped the memory dynamics.
In a follow-up model, long-range recruitment of Martinotti cells produced
surround suppression and promoted the saliency of non-uniform stimuli, and
disinhibition through VIP cells removed that suppression~\cite{krishnamurthy2015longrange}.
Chrysanthidis, Fiebig and Lansner~\cite{chrysanthidis2019introducing} addressed a
different problem. In a spiking BCPNN model with hypercolumns and minicolumns,
the learning rule produces negative as well as positive weights, which a single
pyramidal cell cannot express under Dale's principle. Adding double bouquet
cells, whose axons stay within a minicolumn, turned the learned negative weights
into a plastic disynaptic pathway. Basket cells continued to provide
non-specific feedback inhibition, and the network's learning and memory
function were preserved.

Of the traditions reviewed, this one comes closest to the present thesis, on
two of the three axes. It studies associative memory directly, and it has
already shown that differentiated interneurons matter inside a memory network.
The roles it assigns them are narrow, however. Basket cells implement
competition between minicolumns, Martinotti cells control how long an
attractor lasts and suppress its surround, and double bouquet cells make
learned inhibition compatible with Dale's principle. Where the network does
change mode, as in the sustained activity reported by Fransén and Lansner, the
switch is neuromodulatory. The substrate also differs, since these are spiking
networks whose memories are assemblies of minicolumns selected by soft
winner-take-all competition, not learned SDRs over dendritic segments, and in
the earliest models the stored patterns were wired in rather than learned.
What interneuron-differentiated inhibition could add to this tradition is
therefore not several classes, which it already has, but a division of
encoding and retrieval between them.

\section{Two Operator Mechanisms}
\label{sec:comp-operators}

The preceding sections point to two mechanisms that the HTM tradition lacks and
that the interneuron literature associates with PV and SST cells: a divisive
control of overall response magnitude, and a dendritic control of how feedback
arriving at the apical tuft is combined with feedforward drive. This section
reviews each as a computational operator, states the evidence that links it to
an interneuron class, and states where that link is contested.

\subsection{Divisive Normalisation}
\label{subsec:comp-normalisation}

Carandini and Heeger~\cite{carandini2012normalization} proposed normalisation as
a canonical neural computation. In its standard form, the response of a neuron
is its driving input divided by a constant plus the summed activity of a pool of
neurons,
\begin{equation}
  R_j = \gamma \, \frac{D_j^{\,n}}{\sigma^{n} + \sum_k D_k^{\,n}},
  \label{eq:normalisation}
\end{equation}
where $D_j$ is the drive to neuron $j$, the sum runs over the normalisation pool,
$\sigma$ sets the input level at which normalisation takes effect, and $\gamma$
and $n$ set the gain and the steepness of the response. First developed to
explain responses in primary visual cortex, normalisation has since been
invoked across sensory modalities and for operations such as the effects of
attention, the encoding of value and multisensory integration. Carandini and
Heeger stress that a canonical computation need not have a single
implementation: different circuits, regions or species may compute it with
different components.

For this thesis the contrast with k-winners-take-all is the important point.
Strict kWTA acts on the rank of inputs and imposes an upper bound on the
number of active units. Normalisation acts on their magnitude: each response
is scaled relative to the activity of the pool, and the ordering of inputs is
preserved. Combined with a firing threshold, normalisation can therefore
produce input-dependent recruitment, in which the number of active units
follows the strength and distribution of the input rather than a bound set by
rank. It does not, by itself, silence the population. If every drive in Equation~\ref{eq:normalisation} is multiplied by the same factor, each response approaches a
fixed fraction of $\gamma$ rather than zero, so a proportional increase in input does
not drive the responses toward zero. Whether a normalising operator can
suppress every unit depends on the size of the pool, the distribution of drive
across it and the position of the threshold. In this thesis, complete
suppression is therefore treated as a property of the proposed implementation
to be tested, not as a general consequence of divisive normalisation.
Input-dependent recruitment is the first property examined when a PV-like
normalising operator replaces kWTA in an SDR column.

The attribution of division to PV cells, however, is not settled.
Wilson et al.~\cite{wilson2012division} activated and suppressed PV and
somatostatin-expressing (SOM, here SST) cells optogenetically in mouse visual
cortex and reported a clean dissociation: PV cells mainly divided responses,
changing their gain while preserving stimulus selectivity, whereas SST cells
mainly subtracted from them, shifting response levels and sharpening
selectivity. Studies published at the same time drew different conclusions.
Lee et al.~\cite{lee2012activation} found that activating PV cells sharpened
orientation and direction tuning and improved orientation discrimination in
behaving mice, while activating SST or VIP cells had no significant effect.
Atallah et al.~\cite{atallah2012parvalbumin} found that PV activity transformed
responses linearly, subtracting and scaling them with little change in tuning
width. El-Boustani and Sur~\cite{elboustani2014response} then showed that the
operation performed by SST cells depends on the stimulus: for small stimuli SST
cells respond later than their targets and inhibit subtractively, but for large
stimuli they respond at the same time and inhibit divisively. They concluded
that the arithmetic of inhibition is a dynamic property of the circuit, not a
fixed property of a cell class. Seybold et al.~\cite{seybold2015} reached a
similar conclusion from a different direction, showing that recurrent network
integration can turn subtractive inhibition of single neurons into divisive
inhibition of the population, or the reverse.

This evidence fixes how the operators are used in this thesis. As stated in
Section~\biosec{1.3}{sec:bio-inhibitory}, the synaptic targeting of PV and SST cells is robust, while the
assignment of a fixed arithmetic operation to each class is not. The thesis
therefore treats ``PV as divisive normalisation'' and ``SST as dendritic
gating'' as operator hypotheses: modelling choices motivated by targeting and
by the studies above, whose consequences are tested rather than assumed. This
is also why the comparison with kWTA is framed as a falsification check. If a
PV-like normalising operator in an SDR column behaves no differently from kWTA,
the operator distinction has no computational consequence at this level of
abstraction, and the thesis reports that result.

\subsection{Dendritic and Apical Arbitration}
\label{subsec:comp-apical}

The second operator concerns where inhibition acts on the pyramidal cell rather
than how strongly. Larkum, Zhu and Sakmann~\cite{larkum1999new} showed that layer
5 pyramidal neurons have a second spike-initiation zone in the distal apical
dendrite. Input to the apical tuft alone is strongly attenuated before it
reaches the axon. If a back-propagating action potential from the soma
coincides with distal input within a window of a few milliseconds, however, the
dendrite fires a calcium spike and the neuron emits a burst. Inhibition of the
distal dendrite can block the calcium spike selectively and so prevent the
burst. Larkum~\cite{larkum2013} later proposed this coupling as an organising
principle of cortex. Feedforward input arrives near the soma, feedback from
other areas arrives at the apical tuft in layer 1, and the pyramidal cell itself
acts as an associative element that signals when the two coincide. In this view
the inhibitory control of the dendrite decides whether feedback is allowed to
combine with feedforward drive. Murayama et al.~\cite{murayama2009dendritic}
identified the interneurons responsible in vivo. Martinotti cells in deep
layers, a class of SST interneuron, inhibit the apical tufts of layer 5 neurons
and set the slope and threshold of their dendritic calcium response to sensory
stimulation, and disynaptic inhibition through these cells can abolish
dendritic calcium spikes entirely.

Two models have turned dendritic arbitration into a learning principle.
Urbanczik and Senn~\cite{urbanczik2014learning} proposed a plasticity rule in
which dendritic synapses change so as to reduce the mismatch between the
dendritic potential and the somatic firing of the same neuron. Sacramento et
al.~\cite{sacramento2018dendritic} built on this rule in a multilayer network of
multi-compartment pyramidal neurons. In their model, lateral interneurons that
they identify with SST cells learn to predict the top-down feedback arriving at
the apical dendrite and to cancel it. Whatever feedback is not cancelled is a
local prediction error, and plasticity driven by that error approximates the
backpropagation algorithm without separate learning phases. Together with the
disinhibitory gating of Yang et al.~\cite{yang2016gating}
(Section~\ref{sec:comp-interneuron}), these models establish dendrite-targeting
inhibition as the natural candidate for deciding how much feedback a pyramidal
cell integrates.

Both operators are well developed in their own settings, but neither has yet
been carried into the memory setting this thesis addresses. Normalisation is
defined independently of any circuit, and its attribution to PV cells is
contested. Apical arbitration is established at the level of single cells and
of circuits for gating and learning, but the questions asked of it are
coincidence detection, pathway selection and credit assignment in supervised
learning. In Sacramento et al., the SST-like interneurons cancel predictable
feedback so that errors can drive learning. They do not decide whether the
network stores a new pattern or recalls a stored one, and the model's explicit
aim of avoiding separate phases is the opposite of the mode arbitration
examined here.

\section{Synthesis: Relation to Prior Work}
\label{sec:comp-synthesis}

Table~\ref{tab:comp-prior-work} places the models reviewed in this \unitname{} on the
three axes introduced in Section~\ref{sec:comp-scope}. Read by rows, it shows
that each tradition is strong on one or two axes and silent on the rest. Read by
columns, it shows how each axis is covered.

The HTM tradition and its associative extension supply the substrate. Sparse
distributed representations over minicolumns with dendritic segments, learned
continuously by local rules, are what the present model is built from. Their
inhibition, however, is a single homogeneous k-winners-take-all rule, which
the authors of the lineage themselves describe as a simplification
(Sections~\ref{subsec:comp-htm-inhibition} and~\ref{subsec:comp-naa}). The interneuron-resolved circuit models
supply the opposite. They give PV, SST and VIP cells distinct targets,
dynamics and roles, but in rate-coded or spiking populations without sparse
distributed representations, and in tasks of gating, modulation, error
signalling and physiological fidelity rather than memory. The attractor and
BCPNN tradition comes closest from a third direction. It studies associative
memory in cortical columns and already includes several interneuron classes,
but the roles it assigns them are competition, attractor dwell time, surround
suppression and compliance with Dale's principle. The two operator mechanisms,
divisive normalisation and dendritic arbitration, are each well established,
but outside the memory setting, and the assignment of each to a single
interneuron class remains contested.

These observations are drawn together into the research gap in \gapref{}.

\begin{table}[htbp]
\centering
\scriptsize
\setlength{\tabcolsep}{3pt}
\renewcommand{\arraystretch}{0.88}
\begin{tabular}{>{\raggedright\arraybackslash}p{2.4cm}>{\raggedright\arraybackslash}p{2.4cm}>{\raggedright\arraybackslash}p{2.3cm}>{\raggedright\arraybackslash}p{2.3cm}>{\raggedright\arraybackslash}p{2.1cm}}
\hline
\textbf{Model} & \textbf{Substrate} & \textbf{Inhibition model} & \textbf{Question tested} & \textbf{Encoding/\allowbreak retrieval control} \\
\hline
\multicolumn{5}{l}{\emph{Hippocampal precedent (Section~\ref{sec:comp-scope})}} \\
Cutsuridis et al.~\cite{cutsuridis2010encoding} & Biophysical spiking CA1 microcircuit & Basket, axo-axonic, bistratified, OLM & Storage and recall across theta & External (septal theta) \\
\hline
\multicolumn{5}{l}{\emph{HTM, CLA and SDR tradition (Section~\ref{sec:comp-htm})}} \\
Spatial Pooler, Temporal Memory~\cite{cui2017sp,hawkins2016neurons} & Binary SDRs, minicolumns, dendritic segments & Homogeneous kWTA, global or local & Sparse coding, sequence learning & Absent \\
Iyer et al.~\cite{iyer2022dendrites} & Deep network with dendritic segments & Homogeneous kWTA & Continual multi-task learning & Absent \\
\hline
\multicolumn{5}{l}{\emph{Associative extensions (Section~\ref{sec:comp-assoc})}} \\
SDM~\cite{kanerva1988sdm} & Abstract binary memory & None; radius selection & Storage and noisy recall & External (caller) \\
NAA~\cite{dobric2024naa} & HTM-derived SDRs, multiple areas & kWTA algorithm, no inhibitory cells & Association, SP stability & Absent \\
\hline
\multicolumn{5}{l}{\emph{Interneuron-resolved circuit models (Section~\ref{sec:comp-interneuron})}} \\
Yang et al.~\cite{yang2016gating} & Rate, multi-compartment pyramidal cell & PV, SST, VIP & Pathway-specific gating & Absent \\
Hertäg and Sprekeler~\cite{hertag2019amplifying} & Rate, two-compartment pyramidal cell & PV, SST, VIP & Somato-dendritic redistribution & Circuit-internal, feedback mode only \\
Garcia del Molino et al.~\cite{garciadelmolino2017paradoxical} & Rate populations & PV, SST, VIP & Response reversal & Absent \\
Hertäg and Sprekeler~\cite{hertag2020prediction} & Rate, plastic interneurons & PV, SST, VIP & Prediction-error learning & Absent \\
Keijser and Sprekeler~\cite{keijser2022optimizing} & Spiking, optimised circuit & Emergent PV- and SST-like classes & Compartment-specific inhibition & Absent \\
Billeh et al.~\cite{billeh2020systematic} & Spiking laminar V1, 230{,}000 cells & 17 data-driven cell classes & Physiological fidelity & Absent \\
Moreni et al.~\cite{moreni2025celltype} & Spiking laminar column & PV, SST, VIP per layer & Feedforward vs.\ feedback states & Absent \\
\hline
\multicolumn{5}{l}{\emph{Associative-memory columns (Section~\ref{sec:comp-attractor})}} \\
Fransén and Lansner~\cite{fransen1998horizontal} & Biophysical spiking columns & Basket cells & Completion, rivalry & Neuro\-modulatory \\
Lundqvist et al.~\cite{lundqvist2006attractor} & Spiking hyper- and minicolumns & Basket soft WTA, RSNP cells & Attractor dynamics & Absent \\
Krishnamurthy et al.~\cite{krishnamurthy2012cortical,krishnamurthy2015longrange} & Spiking attractor network & Basket, Martinotti (VIP in 2015) & Dwell time, surround suppression & Absent \\
Chrysanthidis et al.~\cite{chrysanthidis2019introducing} & Spiking BCPNN network & Basket, double bouquet & Learned inhibition, Dale's principle & Absent \\
\hline
\multicolumn{5}{l}{\emph{Operator mechanisms (Section~\ref{sec:comp-operators})}} \\
Normalisation~\cite{carandini2012normalization} & Circuit-independent & Divisive & Canonical computation & Absent \\
Sacramento et al.~\cite{sacramento2018dendritic} & Rate, multi-compartment network & SST-like apical prediction & Supervised credit assignment & Absent \\
\hline
\textbf{This thesis} & \textbf{SDR laminar column} & \textbf{PV-like divisive, SST-like dendritic} & \textbf{Encoding vs.\ retrieval dissociation} & \textbf{External; PV/SST division tested} \\
\hline
\end{tabular}
\caption{Relation of this thesis to prior work on the three review axes:
substrate, inhibition model, and how encoding and retrieval are arbitrated.}
\label{tab:comp-prior-work}
\end{table}